\section{Étude du mouvement d’une luge 4 saisons}

\begin{wrapfigure}{r}{0.4\textwidth}
    \centering
    \includegraphics[width=\linewidth]{ste10_luge_ete_img1.jpg}
\end{wrapfigure}

La luge 4 saisons est une attraction située au pied du mont Mézenc en région Auvergne-Rhône-Alpes. Elle est constituée de luges sur rails qui parcourent un circuit comprenant deux montées et deux descentes.

D’après \textit{https://www.lugikparc}

\subsection{Données}

\begin{figure}[ht]
    \centering
    \includegraphics[width=0.7\linewidth]{ste10_luge_ete_img2.jpg}
\end{figure}

\begin{center}
\begin{tabular}{|l|r|}
\hline
Masse maximale autorisée pour la luge et ses passagers & \SI{150}{\kilo\gram} \\
\hline
Longueur totale du trajet & \SI{1,2}{\kilo\meter} \\
\hline
Durée du trajet & 6 à 7 min \\
\hline
Altitude $z_A$ au point A (départ) & \SI{1365}{\meter} \\
\hline
Altitude $z_B$ au point B & \SI{1401}{\meter} \\
\hline
Altitude $z_C$ au point C & \SI{1375}{\meter} \\
\hline
Vitesse moyenne en montée (luge tirée par un câble) & \SI{7,2}{\kilo\meter\per\hour} \\
\hline
Vitesse maximale sur le circuit (donnée par le constructeur) & \SI{45}{\kilo\meter\per\hour} \\
\hline
Intensité du champ de pesanteur $g$ & \SI{9,81}{\meter\per\second\squared} \\
\hline
\end{tabular}
\end{center}

L’objectif de cet exercice est d’étudier le mouvement du système \{luge + passager\}.
Cette étude est réalisée dans le référentiel terrestre supposé galiléen.

\subsection{Étude du mouvement du système au cours de la première montée (de A vers B)}

La luge est tractée jusqu’au sommet de la montée : elle est tirée par un câble.

On exploite les données obtenues à l’aide de la montre connectée d’un passager pendant la première montée (de A vers B). On repère par des croix la position du système \{luge + passager\} à intervalles de temps constants $\Delta t = \SI{5,0}{\second}$. La trajectoire est donnée ci-après.

\begin{figure}[ht]
    \centering
    \includegraphics[width=0.7\linewidth]{ste10_luge_ete_img3.jpg}
\end{figure}

\medskip
\noindent
\textbf{Q1.} Déterminer la nature du mouvement à l’aide du document précédent. Justifier.

\medskip
\noindent
\textbf{Q2.} Déterminer, en explicitant la démarche utilisée, la valeur de la vitesse moyenne de déplacement du système. Vérifier que ce résultat est en accord avec les données de l’énoncé.

\medskip
Lors de la première montée, les forces qui modélisent les actions mécaniques appliquées au système sont :
\begin{itemize}
    \item la force de traction par le câble $\vec{T}$, supposée constante sur cette montée ;
    \item le poids : $\vec{P}$ ;
    \item la force de frottements : $\vec{f}$ ;
    \item la réaction normale du support sur le système : $\vec{R}$.
\end{itemize}

\medskip
\noindent
\textbf{Q3.} Compléter le schéma du \textbf{document réponse DR2 en page~\AMCpageref{page:luge_ete_DR2}} (à rendre avec la copie), sans souci d’échelle, avec les forces appliquées au système au cours de la montée.

\medskip
\noindent
\textbf{Q4.} Donner une expression du travail de la force de traction $\vec{T}$ appliquée au système lors de la première montée, entre les points A et B.

\medskip
\noindent
\textbf{Q5.} Indiquer, pour chacune des forces appliquées au système, si le travail est moteur, nul ou résistant. Justifier.

\subsection{Étude du mouvement du système lors de la première descente (de B vers C)}

Après la montée, la luge est immobilisée au début de la première descente. Un feu tricolore donne ensuite le top départ pour laisser partir la luge.

La luge n’est pas tractée par un câble pendant la descente.

\medskip
\noindent
\textbf{Q6.} Représenter sur le schéma du \textbf{document réponse DR3 en page~\AMCpageref{page:luge_ete_DR2}} (à rendre avec la copie), sans souci d’échelle, les forces qui modélisent les actions mécaniques appliquées au système au cours de la descente.

On prend l’énergie potentielle de pesanteur $E_{pp} = 0$ pour l’altitude $z = 0$.

\medskip
\noindent
\textbf{Q7.} Établir l’expression de l’énergie mécanique $E_m(B)$ du système au point B en fonction de la masse $m$ du système, la vitesse $v_B$ et l’altitude $z_B$ au point B, et l’intensité du champ de pesanteur $g$.

Donner de même l’expression de l’énergie mécanique $E_m(C)$ du système au point C.

Le passager n’utilise pas le frein de la luge pendant la descente. On considère que la valeur de la force de frottements $\vec{f}$ est négligeable.

\medskip
\noindent
\textbf{Q8.} En déduire, en utilisant le principe de conservation de l’énergie mécanique, que la valeur de la vitesse au point C a pour expression~:
\[
v_C = \sqrt{2 \times g \times (z_B - z_C)}
\]

\textit{Le candidat est invité à prendre des initiatives et à présenter la démarche suivie même si elle n’a pas abouti. La démarche suivie est évaluée et nécessite donc d’être correctement présentée.}

\medskip
\noindent
\textbf{Q9.} Calculer la valeur de la vitesse $v_C$ au point C.

\medskip
\noindent
\textbf{Q10.} La valeur $v_C$ calculée précédemment est supérieure à la valeur maximale de la vitesse pouvant être atteinte par la luge sur le circuit, donnée par le constructeur.
Proposer une explication.

\AMCcleardoublepage
\begin{center}
\textbf{DOCUMENT RÉPONSE \\
À RENDRE OBLIGATOIREMENT AVEC LA COPIE}
\AMClabel{page:luge_ete_DR2}
\end{center}

\textbf{NOM :} \dotfill \textbf{Prénom :} \dotfill

\subsection*{Exercice – Étude du mouvement d’une luge 4 saisons}

\medskip
\textbf{Document réponse DR2 :} Schéma des forces appliquées au système \{luge + passager\} pendant la première montée.

\begin{figure}[ht]
    \centering
    \includegraphics[width=0.6\linewidth]{ste10_luge_ete_img4.jpg}
\end{figure}
\medskip

\textbf{Document réponse DR3 :} Schéma des forces appliquées au système \{luge + passager\} pendant la première descente.

\begin{figure}[ht]
    \centering
    \includegraphics[width=0.6\linewidth]{ste10_luge_ete_img5.jpg}
\end{figure}
\AMCcleardoublepage
